\section{Background}
\label{sec:background}

\paragraph{eBPF as a sensor.}
{\sloppy eBPF lets a user-space program attach small verified programs to kernel
hooks---tracepoints such as \texttt{sched\_process\_exec}, kprobes on
functions such as \texttt{tcp\_connect}, and the raw system-call entry
point---and ship data to user space through a ring buffer~\cite{ebpf-intro}.
The verifier bounds what such a program may do (no unbounded loops, no
unchecked memory access), which is what makes it safe to run in the kernel
and also what makes it a poor place for a detector with state, a model, or
a learned database. Production monitors therefore split the work: eBPF
collects, user space decides~\cite{falco,tetragon,tracee}. \sys keeps that
split and changes only what ``user space decides'' runs in.\par}

\paragraph{WebAssembly as a sandbox.}
A WebAssembly module~\cite{wasm-spec} computes over a single linear memory
that it cannot address outside of, calls the host only through functions
the host chose to import, and has no notion of a file, socket or process
unless an import provides one. WASI~\cite{wasi} is the standard set of such
imports; a host may link all of it, part of it, or an implementation that
grants nothing. Runtimes such as Wasmtime~\cite{wasmtime} compile modules
ahead of time to native code and add two controls the specification does
not: \emph{fuel}, a per-store instruction budget that traps when exhausted,
and \emph{epochs}, a counter the host increments on a timer that traps a
running module when a deadline passes. Together with a resource limiter on
\texttt{memory.grow}, these give a host three knobs---memory, instructions,
wall clock---that it can set per instance and per call. The cost of the
boundary is a copy into linear memory and a native call; prior work has
pushed that transition toward zero~\cite{wasmboxc} and used the same
substrate to isolate browser libraries~\cite{rlbox} and proxy
plugins~\cite{proxy-wasm}. What has not been measured is what it costs
when the code behind the boundary is a host-intrusion-detection rule fed by
the kernel at tens of thousands of events per second; that is the gap
\S\ref{sec:eval} fills.

\paragraph{Syscall-sequence detection.}
The detector \sys ships is deliberately conventional. STIDE~\cite{stide}
records every $n$-gram of syscall numbers seen in normal traces and flags
an $n$-gram it has not seen; it is cheap, needs only normal data, and has
well-known blind spots. Supervised models on windowed features or on the raw
sequence~\cite{creech,kim-lstm} do better on ADFA-LD~\cite{adfa-ld}, the
benchmark we use, at a higher cost. We use STIDE as a cheap first stage
and a small MLP as the scorer, and compare the result to both families in
\S\ref{sec:eval:e1} without claiming it is the best choice.
